\subsection*{Appendix D: The Observation and Artifact Checklist}

When conducting contextual inquiry into live developer or organizational workflows, document:
\begin{itemize}[leftmargin=1.2em,itemsep=0.2ex]
  \item \textbf{Initial task context:} Specific sprint goal, codebase branch, and deployment deadline.
  \item \textbf{Sequence of actions:} Exact command-line tools invoked, browser tabs opened, and documentation consulted.
  \item \textbf{Tooling frictions:} Latency pauses, API rate-limit errors, unhandled exceptions, and manual copy-pasting.
  \item \textbf{Human handoffs:} Slack messages sent to coworkers asking for advice or code reviews.
  \item \textbf{Verification criteria:} How does the practitioner confirm whether an output is correct (manual inspection, unit test, visual check)?
  \item \textbf{Confidentiality protocols:} Redact all API keys, customer identifiers, and proprietary logic immediately.
\end{itemize}

\subsection*{Appendix E: The Opportunity/Problem Evidence Register}

\begin{table*}[t]
\centering
\small
\begin{tabularx}{\textwidth}{>{\RaggedRight\bfseries}p{3.2cm} >{\RaggedRight}X >{\RaggedRight}X >{\RaggedRight}p{3.2cm}}
\toprule
Hypothesis & Supporting Field Evidence & Contradicting Evidence & Next Decisive Test \\
\midrule
Engineers need evaluation feedback, not content & 4 developers report copying prompt templates that fail in production & 6 developers have senior peers who review all pull requests & Paid 1-on-1 pipeline diagnostic \\
\addlinespace
Colleges will fund applied lab modules & 2 department heads express urgent desire for industry curriculum & Academic council approval requires 9 months of advance notice & Pitch a zero-credit winter lab sprint \\
\addlinespace
Enterprise leads need team evaluation clinics & VP reports 6-week delay on core search feature due to hallucinations & Legal refuses to clear vendor without 3 years of audited financials & Run pilot on synthetic public dataset \\
\addlinespace
Working engineers can sustain 6-month study & 12 interviewees state long courses are valuable & Work calendars show only 4 available weekend hours & Offer 2-week sprint; measure dropouts \\
\bottomrule
\end{tabularx}
\caption{The opportunity/problem evidence register template.}
\label{tab:ch6_evidence_register}
\end{table*}

\subsection*{Appendix F: The Acute-Problem Interrogation}

Before committing capital to build an offering, answer these 10 questions in writing:
\begin{enumerate}[leftmargin=1.2em,itemsep=0.2ex]
  \item What specific operational incident occurred in the last 30 days to demonstrate this pain?
  \item What was the measurable economic, operational, or career damage incurred?
  \item How frequently does this precise failure recur across the target segment?
  \item What workaround is currently deployed, and what does it cost in monthly labor?
  \item Who has the authority to approve a commercial replacement check?
  \item What happens to the customer's organization if nothing changes for 6 months?
  \item What good-enough alternative exists if your company never launches?
  \item What observable human behavior proves this is an acute priority?
  \item Can your team deliver a permanent fix within current operational capacity?
  \item What specific negative result would make you abandon this market thesis?
\end{enumerate}

\subsection*{Appendix G: Reconstruct a Purchase as a Timeline}

\begin{verbatim}
Customer Role / Segment:
Context / Project Goal:

T-60d: Background dissatisfaction emerges
T-30d: Acute triggering incident occurs
T-21d: Active search and colleague outreach
T-14d: Evaluation of 2 shortlisted options
T-7d:  Internal debate / Infosec pushback
T-0d:  Contract signed / Payment transacted
T+3d:  Onboarding and initial setup
T+14d: First core value milestone achieved
T+60d: Outcome audit: renewal or churn

Facts Directly Observed:
Facts Customer Reported:
Founder Inferences:
Strongest Alternative Explanation:
Next Critical Experiment:
\end{verbatim}

\subsection*{Appendix H: Empirical Comparison of Two Candidate Segments}

\noindent\textbf{Context:} Bayes Institute conducts preliminary field discovery across 2 distinct hypothetical clusters:

\noindent\textbf{Cluster A (General AI Enthusiasts):} 20 self-identified AI learners recruited from public online forums. 16 express intense enthusiasm for a comprehensive 6-month AI curriculum. However, only 1 has spent more than ₹2,000 on technical books in the past year; 14 report having fewer than 3 hours of free weekend study time; and 12 explicitly state they expect job placement guarantees upon completion.

\noindent\textbf{Cluster B (Production-Assigned Engineers):} 10 practicing backend developers recruited from technical GitHub repositories. 7 recently experienced delayed feature releases due to model hallucination issues; 5 currently spend 6+ hours every weekend building makeshift evaluation scripts; and 4 have corporate access to a ₹50,000 annual individual learning reimbursement budget. When offered a paid, 2-day architectural evaluation diagnostic at ₹15,000, 3 immediately submit code samples for screening.

\noindent\textbf{Strategic Decision:} Disqualify Cluster A immediately: enthusiasm without time, capital, and realistic expectations creates high churn. Focus 100\% of founder discovery on Cluster B: smaller population, but dense with acute pain, budget access, and genuine readiness.
